\section{Stochastic Variational Baum--Welch Convergence}
\label{sec:svb-convergence}

\subsection{Setup}

We consider the idealized infinite-data setting:
each minibatch of $B$ sequence pairs is drawn i.i.d.\ from the
true generative model, so there is no model misspecification.
Stochastic Variational Baum--Welch (SVB) alternates:
\begin{enumerate}
\item \textbf{E-step}: Run forward--backward on a minibatch of $B$ pairs,
  accumulate expected sufficient statistics $\hat T_B$.
\item \textbf{M-step}: Update parameters $\theta_{k+1} = M(\hat T_B)$
  using the closed-form M-step.
\end{enumerate}
At convergence, the true parameters $\theta^*$ are a fixed point of
$\theta \mapsto M\bigl(\expect_{\theta}[T]\bigr)$.

\subsection{BDI sufficient statistics}

For a single sequence pair $(x,y)$ with ancestor length $i$, descendant
length $j$, and evolutionary time $t$, the BDI sufficient statistics are
$T = (\expect[B \mid x,y],\; \expect[D \mid x,y],\; \expect[S \mid x,y])$,
recovered via the score function identity.
The M-step is (Section~\ref{sec:bw-tkf91}):
\begin{equation}\label{eq:mstep}
\hat\insrate = \frac{\bar B}{\bar S + t},
\qquad
\hat\delrate = \frac{\bar D}{\bar S},
\end{equation}
where $\bar B, \bar D, \bar S$ denote averages (or sums) over the minibatch.

\subsection{Variance of Minibatch Estimates}

\subsubsection{Per-pair variance}

Let $\sigma^2_B, \sigma^2_D, \sigma^2_S$ denote the variances of
$\expect[B \mid x,y]$, $\expect[D \mid x,y]$, $\expect[S \mid x,y]$
under the data-generating distribution (randomness over $(x,y)$ pairs
drawn from the model).
For a minibatch of $B$ independent pairs, the sample means satisfy
\[
\Var[\bar B] = \frac{\sigma^2_B}{B}, \qquad
\Var[\bar D] = \frac{\sigma^2_D}{B}, \qquad
\Var[\bar S] = \frac{\sigma^2_S}{B}.
\]

\subsubsection{Parameter variance via the delta method}

By the delta method applied to~\eqref{eq:mstep}, the relative error
of $\hat\insrate$ after one M-step on a minibatch of size $B$ is
\begin{equation}\label{eq:delta-lambda}
\Var[\hat\insrate] / \insrate^2
\;\approx\;
\frac{1}{B}
\left[
  \frac{\sigma^2_B}{\expect[B]^2}
  + \frac{\sigma^2_S}{(\expect[S]+t)^2}
  - \frac{2\,\Cov[B,S]}{\expect[B](\expect[S]+t)}
\right]
\equiv \frac{v_\insrate}{B},
\end{equation}
and similarly for $\hat\delrate$:
\begin{equation}\label{eq:delta-mu}
\Var[\hat\delrate] / \delrate^2
\;\approx\;
\frac{1}{B}
\left[
  \frac{\sigma^2_D}{\expect[D]^2}
  + \frac{\sigma^2_S}{\expect[S]^2}
  - \frac{2\,\Cov[D,S]}{\expect[D]\,\expect[S]}
\right]
\equiv \frac{v_\delrate}{B}.
\end{equation}
The constants $v_\insrate, v_\delrate$ depend on $(\insrate, \delrate, t)$
and the sequence length distribution.
They are the diagonal entries of the inverse Fisher information matrix
(per observation), expressed in the $(\insrate, \delrate)$ parameterization.

\subsubsection{Scaling of per-pair variance with sequence length}

The dominant source of randomness is the ancestor/descendant length pair $(i,j)$.
At stationarity with $\kappa = \insrate/\delrate < 1$:
\begin{itemize}
\item $\expect[i] = \kappa/(1-\kappa)$, \quad
  $\Var[i] = \kappa/(1-\kappa)^2$.
\item $\expect[B \mid i] \approx \insrate\,\expect[S \mid i]$, \quad
  $\expect[S \mid i] \approx i/({\delrate - \insrate})$ for $t \gg 1/(\delrate-\insrate)$.
\item The relative variance $v_\insrate$ is $O(1)$ in the number of links
  (it does \emph{not} decrease as sequences get longer, because the
  M-step aggregates over the pair, not per-link).
\end{itemize}
Thus the per-pair Fisher information is $O(1)$ in sequence length ---
longer sequences do not help much for indel parameter estimation
(in contrast to substitution parameters, where Fisher information
grows linearly with alignment length).

\begin{proposition}[Per-pair Fisher information, long-time limit]
For $\delta = \delrate - \insrate > 0$ and $t \to \infty$:
\[
\expect[B] \to \frac{\kappa}{1-\kappa}, \qquad
\expect[D] \to \frac{\kappa}{1-\kappa}, \qquad
\expect[S] \to \frac{2\kappa}{(1-\kappa)\delta},
\]
and the coefficient of variation of $\hat\insrate$ from a single pair
is $O(1)$ (bounded away from zero).
\end{proposition}

This reflects the fact that, at stationarity, each pair provides a
single ``observation'' of the BDI process endpoint distribution
$P(j \mid i)$ --- the information content is fundamentally per-pair,
not per-residue.

\subsection{Pseudocount representation and its advantages}
\label{sec:pseudocount-view}

The SVB iterates can be represented equivalently in two ways: as an
EMA of data-only sufficient statistics $s_k$ with priors $\alpha$
added inside each M-step, or as an EMA directly in
\emph{posterior pseudocount} space
$\tilde{\alpha}_k = s_k + \alpha$.  In the latter view, the update
rule is
\begin{equation}\label{eq:ema-pseudocount}
  \tilde{\alpha}_k
  \;=\; (1 - \eta_k)\,\tilde{\alpha}_{k-1}
        + \eta_k\bigl(\alpha + (N/|B|)\,s_{\text{batch}_k}\bigr),
\end{equation}
matching~\eqref{eq:svi-update} with the prior $T_0 \equiv \alpha$
absorbed into the EMA target.  By induction, $\tilde{\alpha}_k = s_k + \alpha$
at every step given a matched initial state, so the two formulations
produce identical parameter iterates.

For the Dirichlet-multinomial groups, $\tilde{\alpha}$ is the natural
parameter of the conjugate posterior; for the BDI / GTR groups it is
the augmented sufficient-statistic vector at which the regularised
Q-function is maximised (the underlying priors there are non-conjugate
regularisers\ifdefined\MIXDOMPAPER ; see Section~\ref{sec:bw-mixdom-priors}\fi).  In either
case, the closed-form M-step maximises
\begin{equation}\label{eq:pseudocount-mstep}
  J(\theta \mid \tilde{\alpha}) \;=\; \tilde{\alpha} \cdot \log \theta - Z(\theta),
\end{equation}
where $Z$ is the log-partition function of the family.  For a
Dirichlet-multinomial component (transitions, mixture weights,
class pis, classdist), $J = \sum_i (\tilde{\alpha}_i - 1)\log \theta_i$
up to a constant, with unique maximizer $\theta_i \propto \tilde{\alpha}_i - 1$.
The BDI rates $(\insrate,\delrate)$ follow the same pattern but the
joint $J$ couples them through the geometric fragment length prior
and the M-step is obtained from the quadratic in $\kappa = \insrate/\delrate$
of~\eqref{eq:kappa-quadratic}.  The SVB iteration thus reduces to
updating one pytree of pseudocounts per step, with M-step parameters
computed on demand.

\paragraph{Variance via a single Jacobian.}
Let $f\colon \tilde{\alpha}\mapsto\theta$ denote the M-step map.  The
delta method gives
\begin{equation}\label{eq:delta-pseudocount}
  \Var[\theta_k]
  \;\approx\; J_f(\tilde{\alpha}_k)\,\Var[\tilde{\alpha}_k]\,J_f(\tilde{\alpha}_k)^\top.
\end{equation}
Equations~\eqref{eq:delta-lambda}--\eqref{eq:delta-mu} are the diagonal
entries of~\eqref{eq:delta-pseudocount} in the $(\insrate,\delrate)$
parameterization; in pseudocount form, the same delta method is one
derivative of $f$ away regardless of parameter group.

\paragraph{Polyak--Ruppert averaging is built in.}
Expanding~\eqref{eq:ema-pseudocount},
\[
  \tilde{\alpha}_K \;=\; \sum_{j=1}^K w_{j,K}
    \bigl(\alpha + (N/|B|)\, s_{\text{batch}_j}\bigr),
  \qquad
  w_{j,K} \;=\; \eta_j \prod_{i>j}^K(1-\eta_i).
\]
The pseudocount state is itself a weighted average of batch-level
pseudocounts, and the M-step applied to this average is the
Rao--Blackwell improvement over averaging parameter iterates --- which
matters when $f$ is nonlinear, as for the BDI rate M-step
(\ifdefined\MIXDOMPAPER\S\ref{sec:bw-mixdom}, \fi eq.~\eqref{eq:kappa-quadratic}).  An
effective sample size
\[
  \mathrm{ESS}_K \;=\;
  \frac{\bigl(\sum_j w_{j,K}\bigr)^2}{\sum_j w_{j,K}^2}
\]
quantifies how ``warmed up'' the EMA is and is strictly less than $K$
whenever $\eta_k < 1$.  For a parameter group receiving a
non-negligible contribution in only a fraction $\varepsilon$ of
batches (a rare domain, a rarely-visited fragment state, a
rarely-active class), the per-group effective sample size is
$\varepsilon\,\mathrm{ESS}_K$, not $\mathrm{ESS}_K$ --- exposing the
rare-parameter bottleneck of \S\ref{sec:expected-stats} as a direct
quantity rather than an approximation.

\paragraph{Complete-data Fisher information.}
For a Dirichlet-multinomial component, the Hessian of $J$ at its
maximum is diagonal in natural-parameter coordinates with
$[I_{\text{comp}}(\tilde{\alpha})]_{ii} = 1/\tilde{\alpha}_i$.
Rare categories (small $\tilde{\alpha}_i$) therefore have large
per-component variance, and the bottleneck for each parameter group
can be read off directly as $\tilde{\alpha}_i/\mathrm{ESS}$ ---
a diagnostic in natural-parameter space that replaces the CV-based
convergence recommendation of item~6 in \S\ref{sec:svb-convergence}.

\paragraph{Practical consequence.}
An implementation that carries $\tilde{\alpha}_k$ as its primary state
--- rather than separately tracking $s_k$ and re-adding $\alpha$ in
each M-step --- makes~\eqref{eq:delta-pseudocount}, the
Polyak--Ruppert weights, and the per-group Fisher diagnostics all
computable from the same pytree, with no additional bookkeeping.

\subsection{SVB Convergence Rate}

\subsubsection{Stochastic approximation framework}

SVB is a stochastic fixed-point iteration.
By the results of Capp\'e \& Moulines (2009), under regularity conditions
satisfied by exponential families, the parameter iterates $\theta_k$
satisfy
\begin{equation}\label{eq:clt}
\sqrt{K}(\bar\theta_K - \theta^*) \xrightarrow{d} \mathcal{N}(0, \Sigma)
\end{equation}
where $\bar\theta_K = \frac{1}{K}\sum_{k=1}^K \theta_k$ is the
Polyak--Ruppert average and $\Sigma$ depends on the learning rate schedule.

\paragraph{Direct M-step (no averaging).}
With direct M-step updates (step size 1), the iterate $\theta_K$ has
variance
\begin{equation}\label{eq:nonavg}
\Var[\theta_K] \approx \frac{\Sigma_{\text{per-pair}}}{B}
\end{equation}
after convergence of the transient, where $\Sigma_{\text{per-pair}}$
is the asymptotic covariance of the M-step estimator applied to a
single pair.
This does \emph{not} decrease with $K$ --- the iterates fluctuate
around $\theta^*$ with amplitude $O(1/\sqrt{B})$.

\paragraph{Polyak--Ruppert averaging.}
Averaging the iterates:
\begin{equation}\label{eq:avg}
\Var[\bar\theta_K] \approx \frac{\Sigma_{\text{per-pair}}}{BK}.
\end{equation}
This is optimal: the total number of pairs processed is $N = BK$,
and we achieve variance $\Sigma / N$ regardless of how $N$ is
split between $B$ and $K$.

\paragraph{Exponential moving average (EMA).}
In practice, one often uses $\theta_{k+1} = (1-\eta)\theta_k + \eta\, M(\hat T_B)$
with step size $\eta \in (0,1]$.
The asymptotic variance is
\[
\Var[\theta_\infty] \approx \frac{\eta}{2-\eta} \cdot
  \frac{\Sigma_{\text{per-pair}}}{B}
  \approx \frac{\eta\, \Sigma_{\text{per-pair}}}{2B}
  \quad (\eta \ll 1).
\]
This trades convergence speed for lower asymptotic variance.

\subsubsection{Convergence transient}

The EM rate matrix governs the transient.
For exponential families, the EM convergence rate is
$\rho = 1 - I_{\text{obs}} / I_{\text{comp}}$,
where $I_{\text{obs}}$ and $I_{\text{comp}}$ are the observed and
complete-data Fisher information matrices.
For BDI, the fraction of missing information is typically moderate
($\rho \approx 0.3$--$0.7$ for typical protein evolution times
$t \sim 0.5$--$2$), so the transient dies out in $O(10)$ iterations.

\ifdefined\MIXDOMPAPER
% MixDom-specific subsections moved to a dedicated appendix; see
% Appendix~\ref{sec:svb-convergence-mixdom}.
\fi

\subsection{Practical Recommendations}

\subsubsection{Minibatch size}

For target relative error $\varepsilon$ on indel parameters after $K$ steps
with Polyak--Ruppert averaging:
\begin{equation}\label{eq:minibatch}
BK \geq \frac{v_\theta}{\varepsilon^2},
\end{equation}
where $v_\theta \sim O(1)$ is the per-pair relative variance
from~\eqref{eq:delta-lambda}--\eqref{eq:delta-mu}.

\begin{center}
\begin{tabular}{cccc}
\toprule
Target $\varepsilon$ & $BK$ needed & $B=50, K=?$ & $B=200, K=?$ \\
\midrule
0.10 & $\sim 100\,v_\theta$ & 2 & 1 \\
0.05 & $\sim 400\,v_\theta$ & 8 & 2 \\
0.01 & $\sim 10{,}000\,v_\theta$ & 200 & 50 \\
\bottomrule
\end{tabular}
\end{center}

\paragraph{Rule of thumb.}
Take $v_\theta \approx 5$ as a conservative estimate for typical protein
BDI parameters ($\kappa \approx 0.9$, $t \approx 1$).
Then:
\begin{itemize}
\item $B \geq 50$ ensures that each M-step estimate has
  $\lesssim 30\%$ relative error on indel parameters.
\item $B \geq 200$ brings single-step error below $\sim 15\%$.
\item For rare domains ($w_d \approx 0.05$), multiply $B$ by $1/w_d = 20$.
\item After $K$ steps with averaging, total error scales as $1/\sqrt{BK}$.
\end{itemize}

\subsubsection{Natural gradient}

The BDI complete-data statistics form a curved exponential family
with natural parameters $(\log\insrate, \log\delrate, -(\insrate+\delrate))$
and sufficient statistics $(B, D, S)$.
The natural gradient replaces the M-step update
$\Delta\theta = -I_F^{-1} \nabla_\theta Q(\theta)$
where $I_F$ is the Fisher information matrix of the complete-data family.

For the BDI M-step, the closed-form solution~\eqref{eq:mstep} already
implicitly incorporates the complete-data Fisher information structure.
Stochastic natural gradient (SNatGrad) would precondition the
\emph{stochastic} update direction using $I_F^{-1}$, but since
$\hat\insrate = \bar B / (\bar S + t)$ is already an unbiased estimator
of $\insrate$ (to first order), the benefit of explicit preconditioning
is marginal.

\textbf{Recommendation}: Natural gradient is not worth the implementation
complexity for SVB on TKF-family models.  The closed-form M-step
already provides the natural parameterization.
Use Polyak--Ruppert averaging of the M-step outputs instead.

\subsubsection{Summary of recommendations}

\begin{enumerate}
\item \textbf{Minibatch size}: $B \geq 50$ for general parameters;
  $B \geq 200$ or use averaging if rare domains ($w_d < 0.1$) are present.
\item \textbf{Averaging}: Use Polyak--Ruppert averaging of parameter
  iterates. Total error $\propto 1/\sqrt{BK}$, optimal regardless
  of $B/K$ split.
\item \textbf{Step size}: With direct M-step ($\eta = 1$), iterates
  oscillate with amplitude $O(1/\sqrt{B})$.
  EMA with $\eta = 2/(k+2)$ (Polyak schedule) achieves $O(1/\sqrt{BK})$.
\item \textbf{Decoupled updates}: Substitution parameters converge
  $\sim \bar L$ times faster than indel parameters.
  Consider accumulating indel statistics over $\sim \bar L / B$ minibatches
  before updating.
\item \textbf{Natural gradient}: Not recommended; the closed-form M-step
  is already natural.
\item \textbf{Convergence diagnostic}: Monitor the running coefficient
  of variation of $\hat\insrate_d, \hat\delrate_d$ across minibatches.
  Convergence when CV $< \varepsilon$.
\end{enumerate}

\subsection{Bias diagnostics: Hellinger, ESS, and Fisher readouts}
\label{sec:bias-diagnostics}

The convergence diagnostics above answer ``have the parameter iterates
stabilised?'' but are silent on a distinct question that arises in
practice: when validation log-likelihood \emph{regresses} during
training, is the regression within the noise floor predicted by the
minibatch variance, or is it evidence of a structural problem
(class-assignment drift, over-concentration of a mixture component,
over-fitting of the training likelihood)?  The pseudocount view of
\S\ref{sec:pseudocount-view} makes several readouts available that are
cheap to compute from saved checkpoints and sharply discriminate among
these hypotheses.

\paragraph{Hellinger distance on mixture distributions.}
For any probability vector $p$ (e.g.\ a row of \texttt{classdist[d, f, :]},
a fragment-transition row $\ext^{(d)}_{f,\cdot}$, or a domain-weight
vector), the Hellinger distance
\[
  H(p, q) \;=\; \tfrac{1}{\sqrt{2}}\,\|\sqrt{p} - \sqrt{q}\|_2
  \;\in\; [0, 1]
\]
is the natural metric: bounded, symmetric (unlike KL), and
\emph{locally linear in probability-mass movement}.  A transfer of $\delta$
mass from one category to another produces $H \approx \sqrt{\delta}/\sqrt{2}$,
so a 1\% reassignment registers as $H \approx 0.07$ and a clean
class-swap gives $H = 1$.  The drift hypothesis --- that a mixture
component is discretely reassigned between iterations $k$ and $k{+}1$
and that this reassignment accounts for the val-LL regression --- predicts
one or more iter-to-iter Hellinger values well above the EM-concentration
noise floor, where the noise floor is empirically $\sim 0.005$ on
training runs seeded from an informed mixture profile (e.g.\ IQ-TREE's
C10).  A ``sharp'' drift signature appropriate to a $\sim 10\%$ per-pair
val-LL regression would be $H > 0.15$ at the responsible iter; a flat
trajectory at $H \leq 0.02$ throughout the regression window falsifies
the discrete-drift hypothesis and redirects investigation to continuous
alternatives (over-fitting of the training likelihood, slow limit
cycles, or systematic mis-specification).

\paragraph{Effective sample size per parameter group.}
As derived in \S\ref{sec:pseudocount-view}, the EMA effective sample
size $\mathrm{ESS}_K = (\sum_j w_{j,K})^2 / \sum_j w_{j,K}^2$ is a first-class
quantity in pseudocount space, strictly smaller than $K$ whenever
$\eta_k < 1$.  For a parameter group receiving a non-negligible batch
contribution in only a fraction $\varepsilon$ of minibatches, the
\emph{group-specific} effective sample size is $\varepsilon\,\mathrm{ESS}_K$,
turning the rare-parameter bottleneck of \S\ref{sec:expected-stats} from
a theoretical concern into a direct quantity that can be logged per iter.

\paragraph{Fisher readout of remaining uncertainty.}
Combining the two quantities, the per-component posterior uncertainty
for a Dirichlet-multinomial parameter group is proportional to
$\tilde{\alpha}_i / \mathrm{ESS}$: rare categories with small posterior
pseudocount and small per-group ESS are exactly those whose point
estimates are least certain.  This replaces the CV-of-$\hat\theta$
heuristic of item~6 above with a reading in natural-parameter space
and lets one rank parameter groups by residual uncertainty without
running a second pass of DP.

\paragraph{Suggested dashboard.}
At minimum, log per iter: (i) $\eta_k$ and $\mathrm{ESS}_k$; (ii) for each
mixture group (\texttt{classdist}, \texttt{ext} rows, \texttt{dom\_weights}),
the max Hellinger distance to the previous iter and the entropy of the
current iterate; (iii) $\|\Delta\tilde{\alpha}_k\|_1 / \|\tilde{\alpha}_k\|_1$
per parameter group as an $\tilde{\alpha}$-space convergence diagnostic;
(iv) the Hungarian-matched distance of the current $\mathrm{class\_pis}$ (or
analogous mixture-profile tensor) to the initial informed seed, as a
``drift-from-prior'' signature independent of relabelling.  All are
computable from saved parameters alone; no DP is required.  Taken
together, these give a rapid read-out on whether a val-LL regression
reflects minibatch noise, a discrete drift event, or systematic
over-fitting --- and thus whether the response should be a larger
minibatch, a frozen-component diagnostic, or a stronger prior.

%% ================================================================
%% Maraschino Distillation Fitting Analysis
%% Added 2026-03-29
%% ================================================================

% \subsection{Maraschino Distillation Fitting}
% 
% \subsubsection{Setup}
% 
% The Maraschino pipeline estimates MixDom parameters $\theta$ in three stages:
% \begin{enumerate}
% \item \textbf{Count}: From $N$ pairwise alignments (extracted as cherries
%   from phylogenetic trees), count adjacency transitions
%   $c_{\sigma\sigma'}(a,b,a',b')$ for each transition type
%   $\sigma \to \sigma'$ (where $\sigma \in \{S,M,I,D,E\}$),
%   binned by evolutionary time $\tau$ into $T$ bins.
% \item \textbf{Distill}: Given candidate parameters $\theta$, compute
%   order-1 adjacency frequencies $f_{\sigma\sigma'}(a,b,a',b' \mid \theta, \tau)$
%   via algebraic distillation (Woodbury identity), marginalizing over
%   all latent domain/fragment structure.
% \item \textbf{Fit}: Maximize the multinomial log-likelihood
%   \begin{equation}\label{eq:maraschino-ll}
%     \ell_{\text{Mar}}(\theta) = \sum_{\tau} \sum_{\sigma\sigma'}
%     \sum_{a,b,a',b'} c_{\sigma\sigma'}(a,b,a',b')
%     \log p_{\sigma\sigma'}(a',b' \mid a,b;\, \theta,\tau)
%   \end{equation}
%   where $p_{\sigma\sigma'}$ are the row-normalized distilled frequencies,
%   using Adam gradient ascent.
% \end{enumerate}
% 
% \subsubsection{Bias analysis}
% 
% \paragraph{Source of bias.}
% The distilled frequencies $f_{\sigma\sigma'}(\theta,\tau)$ are the
% \emph{exact} order-1 marginals of the full MixDom Pair HMM --- they
% are obtained by analytic marginalization over the latent
% domain/fragment/class structure via the Woodbury identity, not by
% Monte Carlo or variational approximation.
% Therefore, the predicted conditional probabilities
% $p_{\sigma\sigma'}(a',b' \mid a,b;\, \theta,\tau)$ are the \emph{true}
% order-1 transition probabilities of the MixDom model at parameters $\theta$.
% 
% However, the full MixDom model is \emph{not} an order-1 Markov chain:
% consecutive transitions are correlated through shared latent state
% (which domain we are in, which fragment class, etc.).
% The Maraschino loss~\eqref{eq:maraschino-ll} treats each adjacency
% as an independent draw from the order-1 conditional, ignoring these
% higher-order correlations.
% This is a \textbf{composite likelihood} (specifically, a pairwise
% conditional likelihood), and maximum composite likelihood estimators
% (MCLEs) are consistent but generally inefficient.
% 
% \begin{proposition}[Consistency of Maraschino estimator]
% Let $\theta^*$ be the true MixDom parameters and let
% $p^{\text{true}}_{1}(\cdot \mid \theta^*)$ denote the true order-1
% conditionals.
% Since the Woodbury distillation computes these \emph{exactly},
% $\theta^*$ maximizes $\expect[\ell_{\text{Mar}}(\theta)]$ over $\theta$,
% provided the order-1 marginals identify $\theta$.
% Thus the Maraschino estimator $\hat\theta_{\text{Mar}}$ is
% \textbf{consistent} (converges to $\theta^*$ as $N \to \infty$).
% \end{proposition}
% 
% \paragraph{Identifiability caveat.}
% The order-1 marginals may not distinguish all parameter configurations.
% For example, two domain types with identical $(\insrate_d, \delrate_d, r_d, \pi_d)$
% but different weights $v_d$ produce the same order-1 frequencies as a
% single merged domain.
% In practice, with enough domains ($N_{\text{dom}} \geq 3$) and
% distinct substitution models per domain, the order-1 marginals are
% highly informative and parameter identifiability holds.
% 
% \paragraph{Finite-sample bias.}
% Like any MLE, the Maraschino estimator has $O(1/N)$ finite-sample bias
% from the curvature of the log-likelihood.
% This is negligible when $N \gg p$ (number of parameters).
% For MixDom with $N_{\text{dom}} = 8$ domains, $p \sim 100$ parameters,
% and typical Pfam training sets with $N \sim 10^4$--$10^6$ pairs, the
% bias is negligible relative to the sampling variance.
% 
% \subsubsection{Variance and efficiency}
% 
% \paragraph{Variance of the count tensor.}
% The total adjacency count $C = \sum c_{\sigma\sigma'}(\cdot)$ is
% approximately equal to the total number of alignment columns across
% all $N$ pairs (each column contributes one adjacency to its successor).
% For typical proteins with mean alignment length $\bar{L} \approx 200$,
% $C \approx N \bar{L} \approx 200 N$.
% 
% The variance of each cell $c_{\sigma\sigma'}(a,b,a',b')$ is
% multinomial: $\Var[c] = C \cdot p(1-p)$ where $p$ is the true cell
% probability.  This gives relative error
% $\text{CV}[c] = \sqrt{(1-p)/(Cp)} \approx 1/\sqrt{Cp}$.
% 
% \paragraph{Comparison with the MLE.}
% The full MLE (Baum--Welch) uses the complete-data likelihood,
% exploiting the Markov structure of the latent chain.
% The Maraschino estimator uses only pairwise (order-1) marginals,
% which is a composite likelihood.
% By standard results on composite likelihood inference
% (Varin, Reid \& Firth, 2011):
% 
% \begin{equation}\label{eq:godambe}
%   \Var[\hat\theta_{\text{Mar}}]
%   = H^{-1} J\, H^{-1}
%   \;\geq\; I_{\text{obs}}^{-1}
%   = \Var[\hat\theta_{\text{MLE}}]
% \end{equation}
% where $H = -\expect[\nabla^2 \ell_{\text{Mar}}]$ is the Hessian of the
% composite log-likelihood and $J = \Var[\nabla \ell_{\text{Mar}}]$
% is the variability matrix (which accounts for the correlation between
% adjacent transitions that the composite likelihood ignores).
% The inequality is the Godambe sandwich: the MCLE is never more
% efficient than the full MLE.
% 
% \paragraph{How large is the efficiency loss?}
% The efficiency ratio $e = |I_{\text{obs}}^{-1}| / |H^{-1} J H^{-1}|$
% depends on how much information the higher-order (order $\geq 2$)
% correlations carry about $\theta$.
% For MixDom:
% \begin{itemize}
% \item \textbf{Substitution parameters} ($Q, \pi$): The order-1
%   marginals $f_{MM}(a,b,a',b')$ already capture most of the substitution
%   information (each aligned pair contributes independently).
%   Efficiency loss is minimal ($e \gtrsim 0.9$).
% \item \textbf{Extension rate $r$}: Determines fragment length
%   distribution.  Order-1 frequencies capture the MM$\to$MM self-loop
%   probability, which is the primary signal.  Moderate efficiency
%   ($e \approx 0.7$--$0.9$).
% \item \textbf{Indel rates $(\insrate_d, \delrate_d)$}: These are the
%   hardest.  The transitions M$\to$I, M$\to$D, and I$\to$M carry
%   information, but the \emph{run length} distribution of insertions
%   and deletions (order $\geq 2$) also contains information about
%   $(\insrate_d, \delrate_d)$ that the order-1 model cannot exploit.
%   Efficiency may be lower ($e \approx 0.5$--$0.8$), particularly
%   for the ratio $\kappa_d = \insrate_d/\delrate_d$.
% \item \textbf{Domain weights $v$}: Determined by the mixture structure.
%   Order-1 frequencies encode this well (the marginal emission
%   distribution reflects $v$).  Efficiency is high ($e \gtrsim 0.9$).
% \end{itemize}
% 
% \subsection{Effective sample size: Maraschino vs.\ SVB}
% 
% \paragraph{The counting advantage.}
% Maraschino processes all $N$ training pairs in a single batch
% (the count tensor), while SVB processes $B$ pairs per minibatch.
% After Polyak--Ruppert averaging over $K$ SVB steps,
% the effective sample size of SVB is $BK$ pairs.
% The Maraschino estimator sees all $N$ pairs simultaneously.
% 
% For the variance comparison, using~\eqref{eq:avg}
% and~\eqref{eq:godambe}:
% \begin{equation}\label{eq:variance-ratio}
%   \frac{\Var[\hat\theta_{\text{SVB}}]}{\Var[\hat\theta_{\text{Mar}}]}
%   \;\approx\;
%   \frac{N}{BK} \cdot \frac{1}{e}
% \end{equation}
% where $e \in (0.5, 1.0)$ is the composite-likelihood efficiency factor.
% Even with the efficiency penalty, Maraschino has lower variance
% whenever $N/e > BK$, i.e., whenever it has seen substantially more
% data than SVB.
% 
% \paragraph{Wall-clock advantage.}
% The key computational difference:
% \begin{itemize}
% \item \textbf{Maraschino step}: distill MixDom (Woodbury, $O(N_{\text{dom}} \cdot |\mathcal{A}|^2)$
%   per tau bin), compute loss and gradient.
%   Cost per step: $\sim 0.1$--$1$ second (independent of $N$,
%   since counts are pre-aggregated).
% \item \textbf{SVB step}: forward--backward on $B$ pairs,
%   $O(B \cdot \bar{L}^2 \cdot (5N_{\text{dom}})^2)$ per step
%   (2D Pair HMM DP).
%   Cost per step: $\sim 1$--$60$ seconds depending on $B$ and $\bar{L}$.
% \end{itemize}
% The counting phase (stage A) is $O(N \bar{L})$, done once.
% A typical 1000-step Maraschino fit processes the equivalent of
% $1000 N$ ``pair-observations'' from the order-1 perspective, while a
% 1000-step SVB run processes $1000 B$ pairs through the full model.
% 
% \subsection{Hybrid Maraschino $\to$ SVB pipeline}
% 
% \paragraph{Rationale.}
% \begin{itemize}
% \item Maraschino provides a fast, consistent estimator using all
%   available data, but with some efficiency loss from the composite
%   likelihood approximation.
% \item SVB is the full MLE (no approximation gap), but converges slowly
%   from a random or poor initialization.
% \item Initializing SVB from the Maraschino estimate places it near
%   $\theta^*$, eliminating the transient and requiring only a few
%   iterations to refine.
% \end{itemize}
% 
% \paragraph{Expected convergence.}
% Let $\hat\theta_{\text{Mar}}$ be the Maraschino estimate with error
% $\|\hat\theta_{\text{Mar}} - \theta^*\| = O(1/\sqrt{N})$.
% Starting SVB from $\hat\theta_{\text{Mar}}$:
% \begin{enumerate}
% \item The EM transient (Section~3.2) is essentially eliminated, since
%   $\hat\theta_{\text{Mar}}$ is already within the basin of $\theta^*$.
%   Without warm start, the transient requires $O(|\log \varepsilon| / |\log \rho|)$
%   iterations; with warm start, $O(1)$ iterations suffice.
% \item SVB corrects two sources of suboptimality:
%   \begin{itemize}
%   \item The composite-likelihood efficiency gap (factor $1/e$).
%   \item Any remaining bias from insufficient Adam convergence or
%     finite learning rate in the Maraschino fit.
%   \end{itemize}
% \item After $K_{\text{refine}}$ SVB steps with minibatch $B$, the
%   total variance is
%   \begin{equation}\label{eq:hybrid-var}
%     \Var[\hat\theta_{\text{hybrid}}]
%     \approx I_{\text{obs}}^{-1} / (BK_{\text{refine}}),
%   \end{equation}
%   achieving full MLE efficiency.  The Maraschino initialization
%   ensures $K_{\text{refine}} \ll K_{\text{cold}}$ (the number of
%   steps needed from random initialization).
% \end{enumerate}
% 
% \paragraph{How many SVB refinement steps?}
% The refinement phase needs enough steps for the Polyak average to
% wash out the Maraschino initialization bias.
% Since $\hat\theta_{\text{Mar}}$ has error $O(1/\sqrt{N})$, and each
% SVB step reduces the error by factor $\rho$, we need
% \[
%   K_{\text{refine}} \gtrsim \frac{|\log(1/\sqrt{N})|}{|\log \rho|}
%   \sim \frac{\log N}{2 |\log \rho|}.
% \]
% For $N = 10^5$ and $\rho = 0.5$, this is $K_{\text{refine}} \gtrsim 17$
% steps.  In practice, $K_{\text{refine}} = 50$--$100$ should be ample.
% 
% \subsection{Practical recommendations}
% 
% \begin{center}
% \begin{tabular}{lp{6.5cm}l}
% \toprule
% Method & When to use & Convergence \\
% \midrule
% \textbf{Maraschino only} &
%   Fast exploration; hyperparameter sweeps over ($N_{\text{dom}}$, $N_{\text{frag}}$,
%   $N_{\text{class}}$); initial parameter estimates when many model
%   configurations must be screened. &
%   $O(1/\sqrt{N})$, minutes \\
% \textbf{SVB only} &
%   Small datasets ($N < 1000$) where the full model's higher-order
%   information matters; when the order-1 approximation is suspected
%   to be insufficient. &
%   $O(1/\sqrt{BK})$, hours--days \\
% \textbf{Hybrid} &
%   Production training: Maraschino initializes, SVB refines.
%   Best of both worlds.
%   Required when indel parameters or rare-domain weights must be
%   estimated to high precision. &
%   $O(1/\sqrt{BK})$, faster \\
% \bottomrule
% \end{tabular}
% \end{center}
% 
% \paragraph{Recommended workflow.}
% \begin{enumerate}
% \item \textbf{Stage A} (count): Extract cherries, count adjacencies,
%   bin by $\tau$.  Cost: $O(N \bar{L})$, done once.
% \item \textbf{Stage B} (Maraschino fit): Run Adam for 1000--5000 steps
%   with L-BFGS refinement.  Takes minutes.  Produces
%   $\hat\theta_{\text{Mar}}$.
% \item \textbf{Stage C} (optional SVB refinement): Initialize SVB from
%   $\hat\theta_{\text{Mar}}$.  Run 50--200 steps with $B \geq 50$.
%   Use Polyak averaging.  Achieves full MLE efficiency.
% \item \textbf{Diagnostics}: Compare Maraschino and SVB estimates.
%   If $\|\hat\theta_{\text{SVB}} - \hat\theta_{\text{Mar}}\| \ll
%   \|\hat\theta_{\text{Mar}}\|$, the order-1 approximation is adequate
%   and Stage C can be skipped in future runs with similar data.
% \end{enumerate}
% 
% \paragraph{When is Maraschino alone sufficient?}
% The composite-likelihood efficiency gap is smallest when:
% \begin{itemize}
% \item Domains are well-separated in substitution space (the order-1
%   frequencies cleanly distinguish domains).
% \item Fragment extension rates $r_d$ are moderate ($0.3 \leq r_d \leq 0.8$),
%   so the geometric fragment length distribution is well-characterized
%   by the order-1 self-loop probability.
% \item The number of domains $N_{\text{dom}}$ is small ($\leq 10$),
%   so there are few latent-state correlations to miss.
% \item Training data is abundant ($N \gg 10^4$ pairs), so the
%   composite-likelihood variance is dominated by the $1/\sqrt{N}$
%   scaling rather than the efficiency factor.
% \end{itemize}

%% ================================================================
%% Expected Statistics and Linearized Convergence
%% Added 2026-03-29
%% ================================================================

\section{Expected Statistics and Linearized Convergence}
\label{sec:expected-stats}

The preceding sections characterized SVB convergence in terms of
abstract per-pair variances $\sigma^2_B, \sigma^2_D, \sigma^2_S$.
We now derive these quantities in closed form at stationarity,
giving concrete convergence estimates for the BDI parameter group
of any TKF-family model.\ifdefined\MIXDOMPAPER{}  The corresponding
specialisation to the MixDom hierarchy of top-level, per-domain, and
intra-fragment chains is collected in
Appendix~\ref{sec:svb-convergence-mixdom}.\fi

\subsection{BDI expected statistics at stationarity}

\begin{lemma}[Expected BDI sufficient statistics]
\label{lem:expected-suffstats}
Consider a single BDI process with insertion rate $\insrate$,
deletion rate $\delrate$, $\kappa = \insrate/\delrate < 1$,
$\delta = \delrate - \insrate > 0$,
and define $\alpha = e^{-\delrate\evoltime}$,
$\Phi = 1 - \alpha$.
For an ancestor--descendant pair $(i,j)$ with ancestor length~$i$
drawn from the stationary distribution $i \sim \mathrm{Geom}(\kappa)$
(i.e.\ $P(i) = (1-\kappa)\kappa^i$ for $i = 0,1,2,\ldots$),
the expected BDI sufficient statistics are:
\begin{align}
\expect_\pi[i] &= L \;=\; \frac{\kappa}{1-\kappa}
\label{eq:stat-L}
\\[4pt]
\expect_\pi[S] &= L\,\evoltime
\label{eq:stat-ES}
\\[4pt]
\expect_\pi[B] &= \insrate\,(L + 1)\,\evoltime
  \;=\; \frac{\insrate\,\evoltime}{1 - \kappa}
\label{eq:stat-EB}
\\[4pt]
\expect_\pi[D] &= \delrate\,L\,\evoltime
  \;=\; \frac{\delrate\,\kappa\,\evoltime}{1 - \kappa}
  \;=\; \expect_\pi[B]
\label{eq:stat-ED}
\end{align}
with $M = 1$ (one trajectory endpoint per pair) and $T = \evoltime$
(total observation time for a single process).

The key intermediate result is: for a BDI process starting from
$N(0) = i$ mortal links, the expected number at time~$s$ is
\begin{equation}\label{eq:bdi-mean}
\expect[N(s) \mid N(0) = i]
  = i\,e^{-\delta s} + L\,(1 - e^{-\delta s}),
\end{equation}
so the time-integrated expected count is
\begin{equation}\label{eq:time-integrated}
\expect[S \mid i] = \int_0^\evoltime \expect[N(s) \mid i]\,ds
  = (i - L)\,\frac{1 - e^{-\delta\evoltime}}{\delta} + L\,\evoltime.
\end{equation}
At stationarity, $\expect_\pi[i] = L$, so
$\expect_\pi[S] = L\,\evoltime$.
The birth and death expectations follow from
$\expect[B \mid i] = \insrate\,\expect[S \mid i] + \insrate\,\evoltime$
(the $\insrate\,\evoltime$ term is the immortal link's contribution)
and $\expect[D \mid i] = \delrate\,\expect[S \mid i]$.
The conservation law $\expect_\pi[B] = \expect_\pi[D]$ confirms
that stationarity is maintained.
\end{lemma}

\begin{proof}
Equation~\eqref{eq:bdi-mean} is the standard result for a
linear birth-death process with immigration at rate~$\insrate$
and per-capita death rate~$\delrate$.
The $i$ initial links each survive for $\mathrm{Exp}(\delrate)$ time
(giving the $i\,e^{-\delta s}$ decay), while the immigration
generates new links at rate $\insrate$ from the immortal link plus
the mortal links themselves; the mean converges to the stationary
value $L = \kappa/(1-\kappa)$ as $s \to \infty$.

Integrating~\eqref{eq:bdi-mean}:
\begin{align*}
\expect[S \mid i]
  &= \int_0^\evoltime \bigl[i\,e^{-\delta s}
     + L\,(1 - e^{-\delta s})\bigr]\,ds \\
  &= i\,\frac{1 - e^{-\delta\evoltime}}{\delta}
     + L\Bigl(\evoltime - \frac{1 - e^{-\delta\evoltime}}{\delta}\Bigr) \\
  &= (i - L)\,\frac{1 - e^{-\delta\evoltime}}{\delta}
     + L\,\evoltime.
\end{align*}
Taking $\expect_\pi[i] = L$, the $(i-L)$ term vanishes, giving
$\expect_\pi[S] = L\,\evoltime$.

For births:
$\expect_\pi[B]
  = \insrate\,\expect_\pi[S] + \insrate\,\evoltime
  = \insrate\,(L+1)\,\evoltime
  = \insrate\,\evoltime/(1-\kappa)$.

For deaths:
$\expect_\pi[D]
  = \delrate\,\expect_\pi[S]
  = \delrate\,L\,\evoltime
  = \delrate\,\kappa\,\evoltime/(1-\kappa)
  = \insrate\,\evoltime/(1-\kappa)
  = \expect_\pi[B]$.
\end{proof}

\subsection{Per-pair variance of sufficient statistics}

The randomness in $(B, D, S)$ across pairs arises from the randomness
in the ancestor and descendant lengths $(i, j)$.
At stationarity, $i \sim \mathrm{Geom}(\kappa)$ with
$\Var[i] = \kappa/(1-\kappa)^2$.

Since $\expect[S \mid i]$ is linear in~$i$
(eq.~\eqref{eq:time-integrated}), the variance of $\expect[S \mid i]$
across pairs is:
\begin{equation}\label{eq:var-ES-outer}
\Var_\pi[\expect[S \mid i]]
  = \left(\frac{1 - e^{-\delta\evoltime}}{\delta}\right)^{\!2}
    \cdot \frac{\kappa}{(1-\kappa)^2}
  \;\equiv\; \sigma^2_{S,\text{outer}}.
\end{equation}
Similarly, since $\expect[B \mid i] = \insrate\,\expect[S \mid i] + \insrate\,\evoltime$:
\begin{equation}\label{eq:var-EB-outer}
\Var_\pi[\expect[B \mid i]]
  = \insrate^2\,\sigma^2_{S,\text{outer}},
\end{equation}
and $\expect[D \mid i] = \delrate\,\expect[S \mid i]$:
\begin{equation}\label{eq:var-ED-outer}
\Var_\pi[\expect[D \mid i]]
  = \delrate^2\,\sigma^2_{S,\text{outer}}.
\end{equation}

By the law of total variance,
$\Var_\pi[S] = \Var_\pi[\expect[S \mid i]] + \expect_\pi[\Var[S \mid i]]$.
The inner variance $\Var[S \mid i]$ captures the randomness of the
BDI trajectory given a fixed starting state~$i$.
This is harder to compute in closed form, but for the purpose of
convergence estimates, the outer variance~\eqref{eq:var-ES-outer}
provides a \emph{lower bound} on $\Var[S]$ and is the dominant term
when $\kappa$ is near~1 (long sequences), since both sources of
variance scale as $O(\kappa/(1-\kappa)^2)$.

The covariance structure is also dominated by the outer term:
\begin{equation}\label{eq:cov-BS-outer}
\Cov_\pi[\expect[B \mid i],\, \expect[S \mid i]]
  = \insrate\,\sigma^2_{S,\text{outer}}.
\end{equation}

\subsection{Relative error and per-pair coefficients of variation}

Substituting into the delta-method formulas
\eqref{eq:delta-lambda}--\eqref{eq:delta-mu}
with $\expect[B] = \insrate\,\evoltime/(1-\kappa)$,
$\expect[S] = L\,\evoltime$, and the outer-variance approximation:

Define
\[
\psi = \frac{1 - e^{-\delta\evoltime}}{\delta\,\evoltime}
  \;\in\;(0,1)
\]
(the ratio of the ``regression to mean'' time scale to the total time).
Then:
\begin{align}
\sigma^2_{S,\text{outer}} &= \psi^2\,\evoltime^2\,\frac{\kappa}{(1-\kappa)^2}
\label{eq:sigma-S-psi}
\\[4pt]
\frac{\sigma^2_{S,\text{outer}}}{\expect_\pi[S]^2}
  &= \frac{\psi^2}{\kappa\,/(1-\kappa)^2 \cdot \evoltime^2}
     \cdot \frac{\kappa}{(1-\kappa)^2}
  \;=\; \frac{\psi^2\,(1-\kappa)^2}{\kappa}
  \;\cdot\; \frac{1}{1}
\notag
\end{align}

The relative variance of $\hat\insrate$ from a single pair becomes
(using the outer-variance approximation in~\eqref{eq:delta-lambda},
and noting that the $B$--$S$ covariance
term partially cancels the $B$ and $S$ terms):
\begin{equation}\label{eq:v-lambda-closed}
v_\insrate \;\approx\;
  \frac{\psi^2\,(1-\kappa)^2}{\kappa}
  \left[
    1
    + \frac{\insrate^2}{(\insrate + 1/\evoltime)^2}
    - \frac{2\insrate}{\insrate + 1/\evoltime}
  \right]
  \;=\;
  \frac{\psi^2\,(1-\kappa)^2}{\kappa\,(1+\insrate\evoltime)^2}
\end{equation}
This uses $\expect[S] + \evoltime = (L+1)\,\evoltime = \evoltime/(1-\kappa)$,
so $\insrate/(\insrate + 1/\evoltime) = \insrate\evoltime/(1+\insrate\evoltime) = \kappa\delta\evoltime/(\delta + \delta/(\kappa\cdots))$---more
directly, substitute into~\eqref{eq:delta-lambda}:
\begin{align*}
v_\insrate &= \frac{\insrate^2\sigma^2_{S,\text{outer}}}{\expect[B]^2}
  + \frac{\sigma^2_{S,\text{outer}}}{(\expect[S]+\evoltime)^2}
  - \frac{2\insrate\sigma^2_{S,\text{outer}}}{\expect[B](\expect[S]+\evoltime)} \\
  &= \sigma^2_{S,\text{outer}}
     \left(
       \frac{\insrate^2(1-\kappa)^2}{\insrate^2\evoltime^2}
       + \frac{(1-\kappa)^2}{\evoltime^2}
       - \frac{2\insrate(1-\kappa)^2}{\insrate\evoltime^2}
     \right) \\
  &= \frac{\psi^2\kappa}{(1-\kappa)^2}
     \cdot \frac{(1-\kappa)^2}{\evoltime^2}
     \cdot \bigl(1 + 1 - 2\bigr) \\
  &= 0.
\end{align*}
The outer-variance contribution to $v_\insrate$ vanishes exactly!
This is because $\expect[B \mid i] = \insrate\,(\expect[S \mid i] + \evoltime)$,
so the ratio $\expect[B \mid i]/(\expect[S \mid i] + \evoltime) = \insrate$
is constant in~$i$---the delta-method perturbation is zero.

This means the \textbf{dominant source of variance in $\hat\insrate$
is the inner (trajectory) variance}, not the length distribution.
The per-pair coefficient of variation of $\hat\insrate$ is determined
by the fluctuations in $(B, D, S)$ \emph{given} the endpoint pair
$(i, j)$, which depends on the stochastic path structure of the
BDI process.

\begin{proposition}[Per-pair relative variance from inner fluctuations]
\label{prop:inner-variance}
For $\hat\insrate = \bar B / (\bar S + \evoltime)$ estimated from
a single pair with ancestor length~$i$ drawn from stationarity:
\begin{equation}\label{eq:v-lambda-inner}
v_\insrate \;\approx\;
  \frac{\expect_\pi[\Var[B \mid i, j]]}{\expect_\pi[B]^2}
  + \frac{\expect_\pi[\Var[S \mid i, j]]}{(\expect_\pi[S] + \evoltime)^2}
  - \frac{2\,\expect_\pi[\Cov[B, S \mid i, j]]}
         {\expect_\pi[B]\,(\expect_\pi[S] + \evoltime)}.
\end{equation}
The inner variances $\Var[B \mid i,j]$ and $\Var[S \mid i,j]$
are the posterior variances of the BDI trajectory statistics given
the observed endpoints.  These are computable from the observed
Fisher information (Hessian of the log-likelihood), but do not
simplify to elementary closed forms.

However, we can bound the total per-pair relative variance using
the observed Fisher information.
For a single BDI process observed at one timepoint, the
complete-data Fisher information for $\insrate$ is
$I_{\text{comp}}(\insrate) = \expect[S + \evoltime]/\insrate^2
= \evoltime/(\insrate^2(1-\kappa))$,
and the observed (incomplete-data) information satisfies
$I_{\text{obs}} \leq I_{\text{comp}}$.
The per-pair relative variance is
$v_\insrate \geq 1/(\insrate^2\,I_{\text{obs}}) \geq 1/(\insrate^2\,I_{\text{comp}})$,
giving:
\begin{equation}\label{eq:v-lambda-bound}
v_\insrate \;\geq\; \frac{1-\kappa}{\insrate\,\evoltime}
  \;=\; \frac{1}{\expect_\pi[B]}.
\end{equation}
The per-pair relative variance is at least $1/\expect_\pi[B]$.
When $\expect_\pi[B]$ is small (few births per pair),
the relative error is large.
\end{proposition}

\paragraph{Practical approximation.}
For moderate times ($\delta\evoltime \sim 1$), the missing-information
fraction is $\rho \approx 0.3$--$0.7$, so
$I_{\text{obs}} \approx (1-\rho)\,I_{\text{comp}}$
and $v_\insrate \approx (1-\kappa)/(\insrate\,\evoltime\,(1-\rho))$.
Taking $\rho \approx 0.5$ as a rough estimate:
\begin{equation}\label{eq:v-lambda-approx}
v_\insrate \;\approx\; \frac{2}{\expect_\pi[B]}
  \;=\; \frac{2(1-\kappa)}{\insrate\,\evoltime}.
\end{equation}
Similarly:
\begin{equation}\label{eq:v-mu-approx}
v_\delrate \;\approx\; \frac{2}{\expect_\pi[D]}
  \;=\; \frac{2(1-\kappa)}{\delrate\,\kappa\,\evoltime}.
\end{equation}

\ifdefined\MIXDOMPAPER
% MixDom-specific expected statistics, convergence rate estimates,
% and the discussion of why top-level indel rates are hardest are
% collected in Appendix~\ref{sec:svb-convergence-mixdom}.
\fi
